\honest{My Naturalistic Awakening}{Eliezer Yudkowsky}{2008}

Yesterday my younger self was giving ground to the truth one foot at a time. I quote a saying against this: ``Surrender to the truth as quickly as you can.'' \nb{The saying is from the author's own ``Twelve Virtues of Rationality'' (2006).}

My memory of those years is poor, and I am not sure of the order of events. But if I must pick the moment my folly broke, it is when I first understood, in full generality, the notion of an optimization process. \nb{``The Magnitude of His Own Folly'', five days later, dates the realization to late 2002.}

In 2002 I was writing ``Levels of Organization in General Intelligence''. I thought I was writing about AI, but I was still looking to the human brain for inspiration. I saw natural selection and intelligence as a dichotomy: blind selection against foresight, playing everything out in reality against reasoning by simulation. Yet natural selection made human intelligence, so our brains, though not our thoughts, bear its signature. \nb{The paper bears this out; a footnote calls the comparison a ``loose analogy'' rather than a ``special case.''} I still think the difference is shattering, and it ``drives me up the wall'' when people call both processes ``evolutionary''. They differ almost absolutely in important ways, though some concepts, such as consequentialism and cross-domain generality, describe both. Still, seeing them as a dichotomy showed the limits of my vision, like thinking of fruit as a dichotomy between apples and strawberries. \nb{That lumping is an established position, going back at least to Donald Campbell's ``Blind variation and selective retention in creative thought as in other knowledge processes'' (1960).}

Emil Gilliam pointed out that my view of AI looked much like my view of human intelligence. I did not want to build an AI in the image of a human mind, but I had described levels of organization in human thinking and proposed no different ones for the AI. Gilliam asked whether I hewed too close to the human line. I replied that a ``Completely Alien Mind Design'' would be too hard for human engineers, because we could not understand something so alien while building it. I have heard that excuse many times since. What you understand, you can usually reshape into almost any form, keeping its essence; when you do not understand flight, you suppose a flying machine needs feathers. So I was still starting from the human architecture.

What broke my attachment was ``an embarrassing confession'': an unfinished science-fiction story. Its plot device, something like an Outcome Pump, was a physical effect across time, neither cognitive nor evolutionary, that narrowly constrained the possible outcomes. Because it was ``just a story,'' I was free to work it out logically: C was constrained to happen, so B, in the past, was constrained to happen, so A, which led to B, was too. Two points make a dichotomy, and you imagine them opposed; three force you to generalize. With human intelligence, natural selection and the plot device, I generalized: an optimization process squeezes the future into a narrow region of the possible. Shane Legg's collection of 71 definitions shows that this is less obvious than it seems. \nb{One definition in the collection, Lenat and Feigenbaum's, is close to it: ``the power to rapidly find an adequate solution in what appears a priori (to observers) to be an immense search space.''}

Many definitions by AI researchers speak of ``solving problems'' or ``achieving goals'', and only hindsight makes that the same as squeezing the future. ``A goal is a mentalistic object.'' Electrons have no goals and solve no problems, and imagining a goal means imagining an agent full of wanting. \nb{Cybernetics had given ``goal'' a physical reading in 1943: ``a final condition in which the behaving object reaches a definite correlation in time or in space with respect to another object or event.''} If intelligence is achieving goals, it seems sensible to argue whether some goals are better than others, to speak of the wisdom needed to judge goals, of a system changing its goals, of the free will needed to choose plans, or of an AI realizing that its goals are not what its programmers meant. None of these statements translates when you picture an Outcome Pump.

Seeing through the word ``mind'' to a physical process that, just by obeying the laws of physics, squeezes its future into a narrow region was, ``for me at least'', a naturalistic enlightenment. Confusion drifted away. I saw the work that intelligence does; ``smart'' was no longer a property but an engine, echoing the outer universe in its inner part and so steering it. In the same flash I saw that a mind must give off waste heat to obey thermodynamics.

In 2001 I had treated Friendly AI as a precaution, and I did not know whether a superintelligence could really turn its future into paperclips. \nb{Yudkowsky's \textsc{Creating Friendly AI} (2001) had described an AI whose only supergoal is the Riemann Hypothesis, as one way to get a catastrophe. What its author knew is the memoir's to say.} Now I saw the optimizer at work: sensory information coming in, motor instructions going out, and in between a model linking possible actions to outcomes and a utility function over the outcomes. Put in a utility function, and it would steer the future anywhere.

I saw at last that my 1997 design would have turned its future light cone into ``generic tools'': computers with no programs to run, stored energy with no use. It would have wiped out the human species for nothing. I had missed this ``for six damned years''. I thought: ``I've been stupid.'' \nb{``The Importance of Saying `Oops' '' (2007) set these very words against what Enron's executives said.}
